\documentclass[11pt]{article}
\usepackage[margin=25mm]{geometry}
\usepackage{fontspec}
\setmainfont{Latin Modern Roman}
\usepackage{amsmath,amssymb,mathtools}
\usepackage{unicode-math}
\setmathfont{Latin Modern Math}
\usepackage{xurl,hyperref}
\hypersetup{colorlinks=true,urlcolor=blue}
\setcounter{secnumdepth}{0}
\setlength{\parindent}{0pt}
\setlength{\parskip}{5pt}
\setlength{\emergencystretch}{3em}
\newcommand{\sref}[2]{\hyperlink{source:#1:#2}{[S#2]}}
\newcommand{\cataloguescope}[1]{\paragraph{Recorded target.}#1}
\newcommand{\R}{\mathbb R}
\newcommand{\C}{\mathbb C}
\newcommand{\N}{\mathbb N}
\newcommand{\Z}{\mathbb Z}
\newcommand{\Q}{\mathbb Q}
\newcommand{\F}{\mathbb F}
\newcommand{\eps}{\varepsilon}
\begin{document}
\section{Reflection-Positivity Research Questions}\label{statement}

\subsection{Definitions and mathematical statement}

\paragraph{Shared notation.}$\mathbb N=\{1,2,\ldots\}$; the standard number fields have their usual meanings. All additional conventions are stated locally.


\paragraph{Definitions and conventions.}
In its basic form, reflection positivity concerns a Hilbert space $\mathcal E$, a unitary involution $\theta$ and a subspace $\mathcal E_+$ such that $\langle\theta v,v\rangle\geq0$ for $v\in\mathcal E_+$. Quotienting by the null space produces a physical Hilbert space. Osterwalder--Schrader axioms describe Euclidean correlation functions; Wightman axioms describe relativistic fields. A Borchers triple consists of an operator algebra, translation representation and vacuum with compatible wedge-locality properties. Super versions incorporate a $\mathbb Z_2$ grading.

\paragraph{Target retained from the canonical source catalogue.}
\paragraph{Catalogue wording retained for traceability.}
Resolve the following questions about reflection positivity: describe super versions of massive free fields and reflection-positive Hilbert spaces; determine the role of antiunitary involutions; construct Euclidean realizations of Borchers triples and algebraic quantum field theories; find diffeomorphism-invariant examples and a noncommutative-geometric theory; clarify links with complex measures, stochastic quantization, number theory, string theory, one-dimensional conformal field theory, cluster expansions, modular theory, and the Bessis–Moussa–Villani theorem; remove growth assumptions from the equivalence of the Osterwalder–Schrader and Wightman axioms; and characterize the common structural core of these phenomena.

\paragraph{Source scope and qualification.}
This is the report\textquotesingle s bundled problem list. All requested connections, realizations and removal of growth assumptions are preserved below; the common-core question is programmatic and no unlisted set of axioms is supplied as an alleged original theorem. The full primary list also asks for a Lax--Phillips-type normal form for one-parameter groups whose positive part preserves a standard real subspace, reflection positivity in situations without a Euclidean Hilbert space, and duality for multiple reflection-positive involutions. It asks about multiparticle cluster-expansion spectra beyond a mass gap, and OS--Wightman equivalence without growth estimates or regularity assumptions. These specific directions are retained as additional source context to the condensed catalogue wording. The report explicitly distinguishes deep open problems from background questions in its list.
\subsection{Short English statement}
Explain and extend reflection positivity across its field-theoretic, algebraic, geometric and probabilistic appearances.
\subsection{Research attempt}
\paragraph{Model, date and investigated scope.}
OpenAI GPT-6 Astra (configured model \texttt{gpt-6-astra}), ultra reasoning; generated 11 October 2026. Constructs a Laplace-mixture reflection-positive kernel and its quotient space, with an explicit signed-weight failure.

\paragraph{Sources and literature check.}
The complete supplied definition, including source qualifications and the dated literature assessment, was reviewed. The calculation below is an elementary restricted attempt motivated by [S2]; it does not claim novelty or a fresh literature-status determination.

\paragraph{Context and investigated scope.} The complete supplied definition, including source qualifications and the dated literature assessment, was reviewed. The calculation below is an elementary restricted attempt motivated by [S2]; it does not claim novelty or a fresh literature-status determination.

\paragraph{A common positive-kernel mechanism.} Let $\nu$ be a finite positive measure on $[0,\infty)$ and define $C(t)=\int e^{-m|t|}d\nu(m)$. For positive times $t_i$ and complex coefficients $c_i$, reflection in time produces
\[
 \sum_{i,j}\overline c_i c_j C(t_i+t_j)
 =\int\left|\sum_jc_je^{-mt_j}\right|^2d\nu(m)\geq0.
\]
Thus the reflected kernel is positive, and its null vectors are precisely combinations whose Laplace-polynomial representative vanishes $\nu$-almost everywhere. The quotient pre-Hilbert space embeds isometrically in $L^2(\nu)$ by sending a time symbol to $m\mapsto e^{-mt}$. For a single mass it is one-dimensional; several masses give multiple directions.

\paragraph{Positivity of the measure matters.} For two distinct masses with weights $w_1>0>w_2$, the determinant of the reflected $2\times2$ matrix at distinct positive times is
\[
 w_1w_2\bigl(e^{-m_1t_1-m_2t_2}-e^{-m_1t_2-m_2t_1}\bigr)^2<0.
\]
So a signed or complex spectral weight is not a harmless replacement of the positive one.

\paragraph{Limits of the common core.} This kernel calculation does not construct graded massive fields, settle antiunitary involutions, reconstruct Borchers triples or AQFTs, or remove OS--Wightman growth/regularity assumptions. Diffeomorphism and noncommutative examples, complex measures, stochastic quantization, number/string/conformal links, multiparticle cluster spectra, modular/BMV connections, Lax--Phillips normal forms, absent Euclidean Hilbert spaces and multiple-involution duality remain separate source directions. The first gap is extending this elementary kernel positivity to the required field/domain/reconstruction structures without assuming the omitted regularity.

\paragraph{Verification and adversarial audit.}
Checked reflected quadratic form, quotient null space, two-mass determinant and retained the complete bundled source program. Model self-audit only; no independent human proof checking. Prompts A and B were applied to this canonical record and its qualified source scope. Checked reflected quadratic form, quotient null space, two-mass determinant and retained the complete bundled source program. Retained unresolved status and zero original-target progress; The source field/reconstruction/domain constructions and removal of growth/regularity assumptions are not achieved.

Coordinator checked the Laplace-kernel positivity representation and strict negative determinant with one negative spectral weight. This additional check was model-assisted, not independent human verification.

\paragraph{Tools, process and contributions.}
Python 3 was used for file preparation and structural checks; no numerical experiment or theorem prover was used. Contribution prompt A was followed by prompt B. GPT-6 Astra with ultra reasoning drafted the research and performed its stated self-audit; any additional coordinator contribution is identified in the verification and audit. A GPT-6-based Codex coordinator assembled the manuscript and checked repository integration; its exact variant and reasoning configuration were not exposed. No human mathematical contribution, independent human proof check or formal proof certification is supplied. Complete derivations and the consulted source references are included here; private reasoning transcripts are not published.

\paragraph{Final status and first remaining gap.}
\textbf{UNRESOLVED}. The source field/reconstruction/domain constructions and removal of growth/regularity assumptions are not achieved.

\section{Completion Estimate}
\noindent\textbf{COMPLETION: 0\%}
The calculation is an auxiliary or restricted attempt and establishes no newly verified portion of the full source target. The source field/reconstruction/domain constructions and removal of growth/regularity assumptions are not achieved.

\subsection{Sources}
\begin{itemize}
\item[\hypertarget{source:statement:1}{S1}] ULAM.AI and the UnsolvedMath Contributors, UnsolvedMath, catalogue entry. Catalogue curation under CC BY 4.0; definitions and English exposition added locally.
\url{https://huggingface.co/datasets/ulamai/UnsolvedMath}.
\item[\hypertarget{source:statement:2}{S2}] Mini-Workshop: Reflection Positivity in Representation Theory, Stochastics and Physics (2014), Oberwolfach Reports, DOI 10.4171/owr/2014/55, pp. 3059--3102. Primary report accessed 2026-10-11; relevant PDF pages 39, 40.
\url{https://doi.org/10.4171/owr/2014/55}.
\end{itemize}
\label{end:statement}
\end{document}
